% =====================================================================
% QUESTION 3 — Collections framework task list  (5 marks)
% Demonstrates: \javafile skeleton + plain itemize with per-task marks
% =====================================================================
\question{3}{5}

\noindent You are given a class \texttt{Product} that models a store item.
Inside the \texttt{main} method of \texttt{Inventory}, complete the five
tasks below using the Java Collections Framework. At least one
\texttt{ArrayList} must be used.

\begin{javafile}{Inventory.java}
class Product {
    private String name;
    private double unitPrice;
    private int quantity;

    public Product(String name, double unitPrice, int quantity) {
        this.name = name;
        this.unitPrice = unitPrice;
        this.quantity = quantity;
    }

    public String getName()      { return name; }
    public double getUnitPrice() { return unitPrice; }
    public int    getQuantity()  { return quantity; }

    public double totalPrice() {
        return unitPrice * quantity;
    }

    @Override
    public String toString() {
        return name + " (" + quantity + " @ " + unitPrice + ")";
    }
}

public class Inventory {
    public static void main(String[] args) {
        // ---------- Complete tasks (i) to (v) below ----------
    }
}
\end{javafile}

\noindent\textbf{Tasks:}
\begin{enumerate}[label=(\roman*), itemsep=2pt, topsep=4pt]
    \item Declare an \texttt{ArrayList<Product>} named \texttt{products}. \hfill [1]
    \item Insert three products: \texttt{"Pen", 5.0, 100},
          \texttt{"Notebook", 40.0, 25} and \texttt{"Eraser", 3.0, 200}. \hfill [1]
    \item Sort \texttt{products} in descending order of \texttt{totalPrice()}. \hfill [1]
    \item Remove every product whose quantity is less than \texttt{30}. \hfill [1]
    \item Print each remaining product, then print the total inventory value
          (the sum of all \texttt{totalPrice()} values). \hfill [1]
\end{enumerate}
